%!TeX root=../thesis.tex

\chapter{Research Design}
\label{cap:design-pesquisa}

This chapter describes the empirical approach adopted to evaluate and refine the
Arch-Stage tool. The design follows a mixed-methods strategy aligned with the thesis
objective: understanding how software engineering teams use a tool for managing
architectural hypotheses in a real-world setting, and identifying refinements needed
to improve its practical utility.

% ---------------------------------------------------------------------------
\section{Empirical Context}
\label{sec:contexto-empirico}

% TODO: insert anonymized description of the company and product domain.
The study takes place within \textbf{CompanyX}, an organization that develops and
operates \textbf{ProductY}, a scaled software product serving
% TODO: describe domain, approximate size, and user base.
a significant user base. Multiple software engineering teams collaborate on different
subsystems of ProductY, each responsible for distinct functional areas but sharing
cross-cutting architectural concerns such as
% TODO: list representative quality attributes (e.g., scalability, reliability,
% observability).
scalability and reliability.

\paragraph{Teams and participants.}
% TODO: describe the number of participating teams, team sizes, and experience levels.
At least two teams are involved in the study, representing varying levels of
architectural experience---from teams with a senior architect to teams where
architectural responsibility is distributed informally among developers.

\paragraph{Tool deployment.}
Arch-Stage was deployed within CompanyX's Backstage instance, making it accessible
alongside existing developer portal features. Teams were introduced to the ArchHypo
concepts and the Arch-Stage interface through a short onboarding session before
beginning their use of the tool.
% TODO: describe onboarding duration and materials.

% ---------------------------------------------------------------------------
\section{Research Approach}
\label{sec:abordagem}

% TODO: decide primary evaluation focus:
% (a) interaction/usability and team practices with Arch-Stage, or
% (b) classification and evolution of architectural uncertainties managed via
%     Arch-Stage,
% and adjust research questions and analysis plan accordingly.

The research adopts a \textbf{mixed-methods} approach combining qualitative and
quantitative data sources. This choice is motivated by the exploratory nature of the
study---there are no prior studies of tool-supported ArchHypo use---and by the need to
capture both behavioral patterns (how teams use the tool) and perceptual data (how
teams experience the tool).

The study can be characterized as an \emph{embedded case study} with multiple units of
analysis (the participating teams), following guidelines from Yin and
Runeson et al.\ for empirical software engineering research.
% TODO: add specific methodology references (Yin 2018, Runeson et al. 2012).

The design is informed by the approach taken in the TSE evaluation of
ArchHypo~\cite{silva2024tse}, in which a participant-observer role was used to collect
multi-source evidence: Kanban board history, user stories, incident reports, the
ArchHypo tracking spreadsheet, relevant e-mails, and a questionnaire administered to
all team members. That study demonstrated the value of combining artifact-based and
perceptual evidence for understanding both how a technique is applied and how
participants experience it. The present thesis extends this methodological approach to
the tool-mediated setting, where Arch-Stage itself becomes both a data source (its
stored artifacts) and an object of evaluation (its usability and fit).

Critically, the TSE evaluation used manual spreadsheets as its primary tracking medium.
This thesis evaluates a qualitatively different mode of application: one in which the
tool scaffolds hypothesis management, potentially reducing the cognitive overhead and
process-change friction identified in Findings~8 and~9 of that study.

% ---------------------------------------------------------------------------
\section{Data Sources}
\label{sec:fontes-dados}

Three primary data sources are planned:

\begin{enumerate}
  \item \textbf{Arch-Stage artifacts.} The hypotheses, assessments, and technical plans
  recorded by teams in Arch-Stage constitute the principal data set. For each
  hypothesis, the following attributes are available: statement, quality attribute,
  source type, uncertainty rating, impact rating, technical plan items (type,
  description, target date, status), and lifecycle status changes over time. Unlike the
  spreadsheet-based data of the TSE study, these artifacts are captured in a structured
  database, enabling systematic extraction and longitudinal analysis.

  \item \textbf{Questionnaires.} Short questionnaires administered to team members
  after a defined usage period (e.g., one sprint cycle or one month) to capture
  perceptions of usability, perceived value, and difficulties encountered.
  % TODO: define questionnaire instrument (e.g., adapted SUS, TAM, or custom
  % Likert-scale items).

  \item \textbf{Semi-structured interviews.} Interviews with a subset of participants
  (e.g., team leads, developers who created hypotheses, developers who did not)
  to explore themes emerging from the artifact and questionnaire data in greater
  depth.
  % TODO: define interview guide topics and participant selection criteria.
\end{enumerate}

Optionally, the study may include:
\begin{itemize}
  \item \textbf{Observation of meetings.} Attending architecture discussions or
  sprint-planning sessions where Arch-Stage is referenced, to observe how the tool
  mediates conversation.
  % TODO: decide whether observation is feasible given access constraints.

  \item \textbf{Usage logs.} If available, server-side logs of Arch-Stage interactions
  (page views, edits, time spent) could provide objective usage metrics.
  % TODO: check whether Backstage logging is enabled and data is accessible.
\end{itemize}

% ---------------------------------------------------------------------------
\section{Analysis Plan}
\label{sec:plano-analise}

% TODO: refine this section after finalizing research questions and data sources.

\paragraph{Artifact analysis.}
Hypotheses and technical plans will be coded thematically to identify recurring types
of architectural uncertainty (e.g., technology-fit uncertainties, integration
uncertainties, performance uncertainties) and patterns in how teams formulate and
resolve them. Uncertainty and impact ratings will be analyzed descriptively
(distributions, medians) and, if sample size permits, compared across teams.

\paragraph{Questionnaire analysis.}
Likert-scale responses will be summarized with descriptive statistics. Open-ended
responses will be coded using the same thematic framework applied to the artifact
data.

\paragraph{Interview analysis.}
Interview transcripts will be analyzed using thematic coding, following the six-phase
approach of Braun and Clarke. Codes will be organized into themes addressing each
research question (tool usage patterns, needed refinements, types of uncertainty).

\paragraph{Triangulation.}
Findings from the three data sources will be triangulated to strengthen the validity
of conclusions. Where artifact data and interview data converge, confidence in the
finding increases; where they diverge, the divergence itself becomes a finding
requiring interpretation.

% ---------------------------------------------------------------------------
\section{Ethical Considerations}
\label{sec:etica}

% TODO: describe IRB/ethics committee approval status.
% TODO: describe informed consent process, data anonymization, and data storage.
All participants will be asked for informed consent before data collection begins.
Company and product names will be anonymized in all publications. Individual
participant data (questionnaire responses, interview transcripts) will be stored
securely and accessible only to the research team.

% ---------------------------------------------------------------------------
\section{Alignment with Qualification Feedback}
\label{sec:feedback-quali}

This research design responds to the main feedback received during the author's
qualification examination~\cite{limaQual2024}:

\begin{itemize}
  \item \textbf{Narrowed scope}: the study focuses on specific teams and a concrete
  tool rather than broad surveys of software architecture trends. The qualification
  committee found the earlier scope---encompassing both a tertiary SMS and a
  multivocal review as principal contributions---too wide for a master's thesis.

  \item \textbf{Rejection of the ``paradox'' framing}: the ``paradox of SA
  decision-making'' terminology used in the qualification was judged unproductive by
  the committee, as it overstated the nature of what is in practice a design
  trade-off. This thesis does not employ that framing.

  \item \textbf{Limited research questions}: three focused questions replace the
  earlier, broader set.

  \item \textbf{Concrete artifact}: Arch-Stage serves as the tangible object of
  evaluation, grounding the research in observable behavior and artifacts. The
  qualification's proposed questionnaire-based survey was repositioned as contextual
  motivation rather than a primary contribution.

  \item \textbf{Empirical rigor}: the mixed-methods design, triangulation plan, and
  explicit ethical considerations aim for a level of rigor appropriate for a master's
  thesis in empirical software engineering. By modeling the data collection strategy
  on the TSE study's multi-source approach while extending it to the tool-mediated
  context, the design offers a coherent methodological lineage.
\end{itemize}
